\section{Approach}
\label{sec:approach}

\begin{figure}
  \centering
  \begin{tikzpicture}[font=\footnotesize, node distance=5mm and 7mm,
      box/.style={draw, rounded corners=2pt, align=center, minimum height=9mm, inner sep=3pt},
      dec/.style={draw, diamond, aspect=2, align=center, inner sep=0.5pt},
      res/.style={draw, rounded corners=2pt, align=center, fill=gray!12, minimum height=9mm, inner sep=3pt},
      >={Stealth[length=2mm]}]
    % top row
    \node[box] (in) {pair $(A,B)$\\ domain $D$};
    \node[box, right=of in] (q) {queue of\\ regions};
    \node[box, right=of q] (ai) {interval analysis\\ $\hat A(R)$, $\hat B(R)$};
    \node[dec, right=of ai] (rules) {Rules\\ 0--2b$'$};
    \node[res, right=9mm of rules] (dec1) {region\\ decided};
    \node[res, right=of dec1] (out) {region map,\\ verdict,\\ coverage};
    % bottom row
    \node[dec, below=9mm of rules] (both) {both\\ $\set{\ok}$?};
    \node[box, right=9mm of both] (smt) {Z3 bit-vector\\ return query};
    \node[box, below=9mm of ai] (replay) {replay witness,\\ keep it};
    \node[box, below=9mm of q] (split) {split $R$ at\\ code constants};
    % edges
    \draw[->] (in) -- (q);
    \draw[->] (q) -- node[above]{$R$} (ai);
    \draw[->] (ai) -- (rules);
    \draw[->] (rules) -- node[above]{yes} (dec1);
    \draw[->] (dec1) -- (out);
    \draw[->] (rules) -- node[right]{no} (both);
    \draw[->] (both) -- node[above]{yes} (smt);
    \draw[->] (smt) -- node[right]{UNSAT} (dec1);
    \draw[->] (smt.south) -- ++(0,-0.45) -| node[pos=0.2, below]{SAT} (replay.south);
    \draw[->] (both) -- node[above]{no} (replay.east |- both) ;
    \draw[->] (replay) -- (split);
    \draw[->] (split) -- (q);
  \end{tikzpicture}
  \caption{Overview. The verdict rules decide a region from the interval analyses of both
  versions; only regions where both provably return normally reach the solver. Undecided
  regions, and regions with a replayed witness, are split until a region budget is reached.}
  \label{fig:overview}
\end{figure}

\Cref{fig:overview} gives an overview. We first define what we mean by the same behaviour
(\cref{sec:background}). We then describe the analysis, the verdict rules and the refinement
loop (\cref{sec:formal}). Finally we state their guarantees (\cref{sec:guarantees}).

\subsection{Background}
\label{sec:background}

\paragraph{Observations}
A \emph{pair} is two static methods $A$ and $B$ with the same descriptor. We support methods
with \code{int} parameters, an \code{int} or \code{void} result, and only integer
instructions, branches, \code{goto} and the assertion pattern. Calls, the heap and
\code{try}/\code{catch} make a pair unsupported. An input is a vector
$x \in \mathbb{Z}_{32}^n$. Running $P$ on $x$ yields the \emph{observation}
$\mathit{obs}_P(x) = (o, v)$. Here $o$ is the outcome class (\ok{}, \code{assertion error},
\dbz{}, \code{*} for non-termination, \ldots). The value $v$ is the 32-bit return value if
$o = \ok{}$, and $\bot$ otherwise. We do not observe exception messages, timing or the path
taken. $A$ and $B$ \emph{agree on $x$} iff $\mathit{obs}_A(x) = \mathit{obs}_B(x)$. A
region $R = I_1 \times \dots \times I_n$ is a box of integer intervals. The pair is
\emph{preserved on $R$} if it agrees on every $x \in R$. It is \emph{changed on $R$} if it
disagrees on some $x \in R$. Supported methods are deterministic, so $\mathit{obs}_P$ is a
function.

\paragraph{Interval abstract interpretation}
Abstract interpretation~\citep{cousot1977abstract} runs a program on abstract values that
each describe a set of concrete values. We use the interval lattice over Java \code{int}s,
$\mathcal{I} = \set{[l,u] \mid \mathit{MIN} \le l \le u \le \mathit{MAX}} \cup \set{\bot}$,
ordered by inclusion. The concretisation is $\gamma([l,u]) = \set{v \mid l \le v \le u}$ and
the abstraction is $\alpha(S) = [\min S, \max S]$. This pair is a Galois connection. An
abstract operation $\mathit{op}^\sharp$ is sound if two conditions hold. First,
$\mathit{op}(x,y) \in \gamma(\mathit{op}^\sharp(X,Y))$ for all $x \in \gamma(X)$ and
$y \in \gamma(Y)$ on which $\mathit{op}$ does not throw. Second, its error set contains every
exception that $\mathit{op}$ can raise there. The lattice has height about $2^{32}$, so loops
need a widening operator $\nabla$ to reach a fixpoint quickly.

\paragraph{Symbolic execution and SMT}
Symbolic execution~\citep{king1976symbolic} runs a program on symbolic inputs. For each
feasible path $\pi$ it collects a path condition and a symbolic result. We can compare two
loop-free methods by asking an SMT solver whether some input satisfies a path condition of
each version and makes the results differ. Java \code{int} arithmetic wraps, so we use 32-bit
bit-vectors. Their \code{bvsdiv} and \code{bvsrem} match \code{idiv} and \code{irem},
including $\mathit{MIN} / {-1} = \mathit{MIN}$. Unbounded integers give wrong answers near the
limits.

\subsection{Formal Description}
\label{sec:formal}

\paragraph{Region analysis}
For a version $P$ and a non-empty region $R$, the analysis computes an outcome set
$\hat P(R)$ and a return interval $\hat r_P(R)$. It is a forward worklist interpreter
over bytecode offsets. The abstract state at each offset is a tuple of intervals for the
locals and a stack of values. The initial state maps parameter $i$ to $I_i$. The analysis has
four parts.
\begin{itemize}
  \item \emph{Arithmetic.} $+$, $-$ and $\times$ compute the mathematical bounds and return
        $\top$ if a bound leaves the \code{int} range. Negation works the same way. For
        $X / Y$ and $X \,\%\, Y$ we add \dbz{} iff $0 \in Y$. We compute the result
        separately on the negative and the positive part of $Y \setminus \set{0}$, and we join
        the two results (\cref{thm:intervals}).
  \item \emph{Branch refinement.} We follow a conditional branch on $X \bowtie Y$ on each
        side, with $X$ and $Y$ narrowed by the condition or its negation. An empty interval
        marks the side as infeasible. Narrowing a stack value is not enough. Every stack value
        therefore carries a tag $\mathit{src}$ that names the local it was loaded from. The
        tag is cleared when that local is written, and we narrow the local itself too. This is
        how $b \neq 0$ in \code{if (b == 0) return 0} reaches the later division. Intervals
        cannot express a hole, so $\neq$ only removes an endpoint.
  \item \emph{Loops.} At a backward edge we widen the new state. An interval bound that grew
        jumps to $\mathit{MIN}$ or $\mathit{MAX}$. Every reachable backward edge adds the
        outcome \code{*} (``may not terminate''), unless the loop matches the counting pattern
        that \code{javac} emits for \code{for (i = ..; i < n; i++)}. That pattern has a single
        exit test \code{i >= n}, a straight-line body with exactly one \code{i += 1}, and no
        writes to \code{i} or \code{n}. Such a loop always exits. The code also accepts the
        mirrored decreasing counter, but we prove only the increasing case.
  \item \emph{Unsupported code.} Any other instruction makes the result \emph{unsupported},
        with the full outcome set and $\top$ returns. No rule can then fire.
\end{itemize}

\paragraph{Verdict rules}
For a region $R$ we try the rules below in order. The first rule that matches decides.
\begin{center}
\small
\begin{tabular}{@{}llll@{}} \toprule
\# & Condition & Verdict & Scope \\ \midrule
0  & $\hat A(R) = \emptyset$ or $\hat B(R) = \emptyset$ & \unk{} & analysis failure \\
1  & $\hat A(R) = \hat B(R) = \set{e}$, $e \neq \ok$ & \pres{} & all of $R$ \\
2  & $\hat A(R) \cap \hat B(R) = \emptyset$ & \chg{} & all of $R$ \\
2b & both $\set{\ok}$, $\hat r_A(R) \cap \hat r_B(R) = \emptyset$ & \chg{} & all of $R$ \\
2b$'$ & both $\set{\ok}$, $\hat r_A(R) = \hat r_B(R) = [c,c]$ & \pres{} & all of $R$ \\
3  & both $\set{\ok}$ and $\Phi(R)$ UNSAT & \pres{} & all of $R$ \\
3$'$ & both $\set{\ok}$ and $\Phi(R)$ has a model $w$ that replays & \chg{} & witness $w$ only \\
4  & otherwise & \unk{}, split $R$ & \\ \bottomrule
\end{tabular}
\end{center}
Rule~1 needs a \emph{singleton} set. Equal sets with two outcomes prove nothing. In
\cref{lst:example} both sets are $\set{\ok, \dbz}$ on $D$, yet the versions differ at
$(0,0)$. In general $A$ and $B$ may swap their outcomes between inputs.

\paragraph{Return comparison}
For rules 3 and 3$'$ we execute both versions symbolically over \emph{shared} bit-vector
inputs $\vec x$. We collect their normally returning paths $(\pi_{P,i}, r_{P,i})$. Every
division on such a path adds the condition ``divisor $\neq 0$'' to $\pi$. SMT-LIB defines
$x / 0$, while Java throws. Then
\[
  \Phi(R) \;=\; \vec x \in R \;\wedge\; \bigvee_{i,j}\bigl(\pi_{A,i} \wedge \pi_{B,j}
  \wedge r_{A,i} \neq r_{B,j}\bigr).
\]
Before we trust an UNSAT answer, we ask Z3 whether the collected paths cover all of $R$ for
both versions. If they do not, the answer is \unk{}. We replay a model $w$ on an independent
concrete interpreter with Java \code{int} semantics. We report it only if the two observations
really differ. A solver timeout, a backward jump or more than 256 paths also give \unk{},
never \pres{}.

\paragraph{Refinement}
\code{check\_pair} keeps a queue of regions, starting with $D$. We split an undecided region
with \code{split\_region}. We take the first parameter whose interval strictly contains a
constant $c$ pushed by either version (or $0$). The split produces $[l, c-1]$, $[c,c]$ and
$[c+1, u]$. Without such a constant we bisect the widest interval. The children are disjoint
and cover the parent. A region decided by Rule~3$'$ keeps its witness but is still split, so
we can certify the rest of it. A budget of about 200 regions bounds the loop. Regions left over
stay \unk{}. The pair verdict is \chg{} if any witness exists, \pres{} if \pres{} regions
cover $D$, and \unk{} otherwise. The \emph{certified coverage} is the fraction of $D$ in
\pres{} regions, whole-region \chg{} regions and witness points.

\subsection{Theoretical Guarantees}
\label{sec:guarantees}

The checker rests on two abstractions. The \emph{interval abstraction} abstracts sets of
\code{int} values. The \emph{region abstraction} abstracts the observation function of a
version on a box. It maps $\mathit{obs}_P$ restricted to $R$ to the pair
$(\hat P(R), \hat r_P(R))$. Its concretisation is
$\gamma(O, r) = \set{(o, v) \mid o \in O,\; o = \ok \Rightarrow v \in \gamma(r)}$. We prove
both abstractions sound and prove that widening terminates. We then derive the correctness of
the verdicts.

\begin{theorem}[Interval transfer functions]
\label{thm:intervals}
Each abstract operation $+^\sharp, -^\sharp, \times^\sharp, /^\sharp, \%^\sharp$ and
$-^\sharp$ (negation) satisfies the soundness condition of \cref{sec:background}, and
$/^\sharp$ and $\%^\sharp$ report \dbz{} iff $0 \in Y$.
\end{theorem}
\begin{proof}
Let $X = [a,b]$ and $Y = [c,d]$. For $+$, $-$ and $\times$, the bounds of the mathematical
result over the box $X \times Y$ lie at its corners. The operations $+$ and $-$ are monotone in
each argument, and $\times$ is linear in each argument. If both bounds lie in the \code{int}
range, every Java result equals the mathematical one. Otherwise the result is $\top$, which
contains every wrapped value. Negation is the case $[-b,-a]$, and only $-\mathit{MIN}$ wraps.
For division, \code{idiv} throws iff $y = 0$. So ``\dbz{} iff $0 \in Y$'' is exact in both
directions. Only $y \in Y \setminus \set{0}$ reaches the quotient. This set is the union of a
negative and a positive interval, so it suffices to bound each part and join. On a part with
fixed sign, $q(x,y) = \mathit{trunc}(x/y)$ is monotone in $x$ for fixed $y$. For fixed $x$ it
attains its extreme over $y$ at an endpoint. The four corners therefore bound it. Since
$|q| \le |x|$, only $\mathit{MIN}/{-1}$ can overflow, and that case returns $\top$. For
remainder, Java guarantees three facts. The result $x \,\%\, y$ is $0$ or has the sign of $x$.
We have $|x \,\%\, y| \le |y| - 1$. We also have $|x \,\%\, y| \le |x|$. With
$m = \max(|c|,|d|) - 1$ on a part, these facts give the bounds
$[\,0 \text{ or } \max(a,-m),\; 0 \text{ or } \min(b,m)\,]$ that we compute.
\end{proof}

\begin{lemma}[Branch refinement]
\label{lem:refine}
For every $x \in X$, $y \in Y$ with $(x \bowtie y) = h$, the refined intervals
$(X', Y')$ contain $x$ and $y$.
\end{lemma}
\begin{proof}
The case $h = \mathit{false}$ uses the negated condition. On integers this is again one of
$=, \neq, <, \le, >, \ge$. For each condition the bounds follow directly. For example,
$x < y \le d$ implies $x \le d - 1$, and $y > x \ge a$ implies $y \ge a + 1$. For $=$ both
values lie in $X \cap Y$. For $\neq$ we remove a value $v$ from $X$ only when $Y = [v,v]$ and
$v$ is an endpoint of $X$. Narrowing the local named by the $\mathit{src}$ tag is sound. A load
sets the tag, and every write to that local clears it. The tag therefore holds only while the
stack value equals the current content of the local.
\end{proof}

\begin{theorem}[Widening terminates]
\label{thm:widening}
$X \sqsubseteq X \nabla Y$ and $Y \sqsubseteq X \nabla Y$, and every chain
$X_0, X_1 = X_0 \nabla Y_1, X_2 = X_1 \nabla Y_2, \ldots$ is eventually constant.
\end{theorem}
\begin{proof}
$X \nabla Y$ keeps each bound of $X$ that $Y$ does not exceed. It moves the other bound to
$\mathit{MIN}$ or $\mathit{MAX}$. The result therefore contains both arguments. Along a chain
each of the two bounds changes at most once, from its current value to its limit, and it
cannot grow after that. A chain therefore changes at most twice. With $k$ locals and stack
slots at a loop head, the abstract state there changes at most $2k$ times. The worklist
revisits a program point only when its state grows, so it terminates.
\end{proof}

\begin{theorem}[Soundness of the region analysis]
\label{thm:region}
For every supported $P$, non-empty region $R$ and $x \in R$:
$\mathit{out}_P(x) \in \hat P(R)$, and $\mathit{ret}_P(x) \in \gamma(\hat r_P(R))$ if
$\mathit{out}_P(x) = \ok$. We call this property (S).
\end{theorem}
\begin{proof}
We show an invariant. Every concrete state $(\mathit{pc}, \mathit{locals}, \mathit{stack})$
reached from an input in $R$ lies in the concretisation of the abstract state stored at
$\mathit{pc}$, interval by interval. The invariant holds initially, because parameter $i$
starts in $I_i$. Each instruction preserves it by \cref{thm:intervals} and \cref{lem:refine}.
Joins and widenings only enlarge states (\cref{thm:widening}). At the fixpoint the invariant
covers every reachable state. Each concrete step that ends the method is matched by an
abstract step. A \code{return} adds \ok{} and joins the returned interval into $\hat r_P$. A
division adds \dbz{} whenever its divisor may be $0$. Constructing an \code{AssertionError}
adds \code{assertion error}. We reject methods with exception handlers, so every exception
escapes. A non-terminating execution must follow some backward edge infinitely often. The
analysis adds \code{*} for every reachable backward edge except a counting loop. Such a loop
runs its body only while $i < n \le \mathit{MAX}$. So $i + 1$ cannot wrap, and $n - i$
decreases by one in every iteration. We reject the pattern \code{i <= n}, because with
$n = \mathit{MAX}$ it never exits.
\end{proof}

\paragraph{Soundness of the verdicts}
With (S), the verdict rules are correct.
\begin{itemize}
  \item \emph{Rule 1.} By (S), both versions end in $e$ on every $x \in R$. Error observations
        carry no value, so $\mathit{obs}_A(x) = (e,\bot) = \mathit{obs}_B(x)$.
  \item \emph{Rule 2.} By (S), $\mathit{out}_A(x) \in \hat A(R)$ and
        $\mathit{out}_B(x) \in \hat B(R)$. The sets are disjoint, so the outcomes differ for
        every $x$. Rule~0 must run first, because an empty set would make Rule~2 fire
        vacuously.
  \item \emph{Rules 2b and 2b$'$.} These are the same arguments applied to the return
        intervals. We treat an empty return interval like Rule~0.
  \item \emph{Rule 3.} Suppose the versions disagree on some $x \in R$. Both return normally,
        so $\mathit{ret}_A(x) \neq \mathit{ret}_B(x)$. The paths are complete, which we check,
        and execution is deterministic. So $x$ satisfies exactly one $\pi_{A,i}$ and one
        $\pi_{B,j}$, and then $x$ is a model of $\Phi(R)$. This contradicts UNSAT. The proof
        needs the bit-vector encoding, the divisor conditions and the shared inputs.
  \item \emph{Rule 3$'$.} A replayed witness is a concrete disagreement.
\end{itemize}
Refinement terminates, because every split strictly shrinks a finite region and the budget
bounds the number of regions. The children partition their parent, so coverage counts each
input once. Hence a \pres{} verdict is a proof over all of $D$. A \chg{} verdict always comes
with a witness, either a replayed solver model or a point of a region that differs everywhere.
The checker is not complete. Intervals cannot relate two variables, so non-constant return
values need the solver. Loops with loop-dependent return values stay \unk{}. Besides the
proofs, we checked every interval operation with property-based tests and with brute-force
checks on intervals with endpoints near $0$, $\mathit{MIN}$ and $\mathit{MAX}$.
